\scanpage{68}
\fbox{(\uline{Dém} de Matsusaka)}

% slanted
\margin{finitude de $\underline{\mathrm{Pic}}^{\tau}_{X/S}$}
Soit $X$ une surface projective non singulière sur un corps algébriquement clos $k$. \inserted{Soit $H$ un diviseur ample.} Rappelons la formule de R-R. faible
\[
\chi(D)=\frac12D(D-K)+c
\]
où $c$ est une constante. \deleted{Soit $D'$ numériquement équivalent à $0$, on trouve} \inserted{Comme}
\[
\chi(D)=\ell(D)+\ell(K-D)-\dim H^1(\unsure{x|X},\unsure{L|\mathcal{L}}(D))\qquad\text{d'où}
\]
\[
\begin{aligned}
&\deleted{\chi(}\\
&\deleted{\chi(D)}\\
&\ell(D)+\ell\unsure{(K-D))|(K-D)}\geq\frac12D(D-K)+c
\end{aligned}
\]

\deleted{Donc on en conclut}

Supposons p.ex. que $\ell(K-D)=0$, ce qui est le cas si $H.(K-D)\leq0$ i.e. si
\[
\text{a)}\qquad \deleted{\unsure{HK<DH|H(K-D)}} (D-K).H\geq0
\]
Alors on a $\ell(D)\geq\frac12D(D-K)+c$, donc a) et \deleted{\unsure{si on}}
\[
\text{b)}\qquad \frac12D(D-K)+c\geq0
\]
implique $\ell(D)>0$. Les conditions a) et b) sont satisfaites p.ex. pour \deleted{\illegible{1}} $D=nH$, $n$ assez grand. Si alors $\Delta$ \deleted{\unsure{re}} est un diviseur numériquement équivalent à $0$, les conditions a) et b) sont vérifiées pour $\Delta$, donc on a
\[
\ell(D+\Delta)>0
\]
donc $D+\Delta$ est représenté par un diviseur positif. Donc :
% archive sheet number: 67

